\documentclass[11pt]{article}

\usepackage[utf8]{inputenc}         
\usepackage[english]{babel}          

\usepackage{graphicx}               
\usepackage{tabularx}               
\usepackage{multirow}               
\usepackage{url}                

\usepackage[small,bf]{caption2}     
\usepackage{parskip}
\usepackage{titlesec}
\usepackage{amsmath,amssymb}                
\usepackage{overpic,subfigure,tikz}

\input{scidefs}
\begin{document}

\title{Low-Frequency Limit of Wave Field Synthesis}
\date{\today}  
\maketitle

\section{High-Frequency Assumptions in Wave Field Synthesis}

Wave Field Synthesis (WFS) is a well‐established method for reconstructing a desired sound field by driving an array of densely spaced loudspeakers. In three dimensions, WFS is based on the Kirchhoff–Helmholtz integral (KHIE) formulation, which expresses the virtual (target) sound field in terms of a surface integral over a secondary source distribution (SSD). For the special case of an infinite planar SSD, the KHIE can be transformed into a single–layer potential in which the driving functions appear implicitly (given directly by the gradient of the virtual field). For arbitrary (curved) SSD surfaces, the Kirchhoff approximation is applied under the assumption that locally the SSD elements and the wavefronts are approximately planar—a condition valid in the high-frequency regime.

Assume a target field of the form
\begin{equation}
    P(\vx,\omega) = A^P(\vx)\,\te^{\ti\Phi^P(\vx)},
\end{equation}
where \(A^P(\vx)\) is the amplitude and \(\Phi^P(\vx)\) the phase. Then, the integral formulation (for both the exact and approximate cases) is given by
\begin{equation}
    P(\vx,\omega) \approx \int_S \!\left[-2\,\frac{\partial}{\partial n}P(\vxo,\omega)\right] G(\vx-\vxo,\omega)\,\td S(\vxo),
\end{equation}
with \(\frac{\partial}{\partial n}P(\vxo,\omega)\) denoting the normal derivative (at \(\vxo\)) and 
\begin{equation}
    G(\vx-\vxo,\omega) = A^G(\vx-\vxo)\,\te^{-\ti\Phi^G(\vx-\vxo)}
    = \frac{1}{4\pi}\frac{\te^{-\ti k|\vx-\vxo|}}{|\vx-\vxo|}
\end{equation}
being the 3D free–field Green’s function (with \(k=\omega/c\)). In the case of an infinite planar SSD (Rayleigh I integral) the target field is exactly reproduced, whereas for curved SSDs the approximation holds under the condition \cite[Eq.~(2.7.12)]{Blenstein1975}
\begin{equation}
    k\,\rho_{1,2} \gg 1,
\end{equation}
with \(\rho_{1,2}\) being the local principal radii of curvature.

A common simplification is to approximate the gradient of the target field by
\begin{equation}
    \nabla_{\vx}P(\vx,\omega) 
    = \Bigl(\frac{\nabla_{\vx}A^P(\vx)}{A^P(\vx)} + \ti\,\nabla_{\vx}\Phi^P(\vx)\Bigr) P(\vx,\omega)
    \approx \ti\,\nabla_{\vx}\Phi^P(\vx)\,P(\vx,\omega).
\end{equation}
Since in a homogeneous medium we have \(|\nabla_{\vx}\Phi^P(\vx)|=k\), one may write
\[
\nabla_{\vx}\Phi^P(\vx)= k\,\hat{\vk}^P(\vx),
\]
where \(\hat{\vk}^P(\vx)\) is the unit propagation direction. Thus, the high–frequency Kirchhoff approximation becomes
\begin{equation}
    P(\vx,\omega) \approx -2\int_S \underbrace{\ti\,k^P_{\mathrm{n}}(\vxo)P(\vxo,\omega)}_{D_{3D}(\vxo,\omega)} G(\vx-\vxo,\omega)\,\td S(\vxo),
\end{equation}
with \(k^P_{\mathrm{n}}(\vxo)\) denoting the normal component of the local wavevector at the SSD.

In practical scenarios the SSD is implemented as a (linear or curved) contour of loudspeakers, leading to 2.5D Wave Field Synthesis. The derivation of 2.5D WFS is based on the \emph{stationary phase approximation} (SPA). For a fixed receiver position \(\vx\) the SPA states that the main contribution to the integral comes from the SSD location where the phase is stationary. In the 2.5D approach the SPA is applied first in the vertical direction (reducing a surface integral to a line integral) and then, if necessary, also in the horizontal direction. This series of approximations is valid only in the high–frequency region.

To summarize, the following high–frequency assumptions are employed in 2.5D WFS:
\begin{itemize}
    \item The classical formulation approximates the required normal derivative under high–frequency conditions.
    \item The 2.5D Kirchhoff approximation uses the SPA for dimensionality reduction.
    \item For non–planar SSDs the Kirchhoff approximation itself is only valid at high frequencies.
\end{itemize}
Violating these conditions typically results in a significant drop in the synthesized sound field’s amplitude below a cutoff frequency that depends on both the receiver position and the target field model. In the following we focus on the validity of the SPA, which imposes the most stringent high–frequency requirement.

\section{Validity of the Stationary Phase Approximation}

To isolate the effect of the SPA, consider the Rayleigh integral over an infinite \(xz\)–plane:
\begin{multline}
    P(\vx,\omega) = -\iint_{-\infty}^{\infty} 2\,\frac{\partial}{\partial y}P(\vxo,\omega)\,G(\vx-\vxo,\omega)\,\td x_0\,\td z_0 \\
    \approx \iint_{-\infty}^{\infty} 2\,k_y(\vxo)P(\vxo,\omega)\,G(\vx-\vxo,\omega)\,\td x_0\,\td z_0,
    \label{eq:Rayleigh}
\end{multline}
or, in polar form,
\begin{multline}
    P(\vx,\omega) \approx 2\iint_{-\infty}^{\infty} \underbrace{\ti\,k_y(\vxo)A^P(\vxo)A^G(\vx-\vxo)}_{A(\vx,\vxo)} \, \te^{-\ti\underbrace{\Bigl[\Phi^P(\vxo)+\Phi^G(\vx-\vxo)\Bigr]}_{\Phi(\vx,\vxo)}} \\
    \times\,\td x_0\,\td z_0.
    \label{eq:RayleighPolar}
\end{multline}
Our goal is to reduce the two–dimensional surface integral (over \(x_0\) and \(z_0\)) to a one–dimensional line integral along \(x_0\) by applying the SPA in the vertical (\(z_0\)) direction.

\subsection{General Formulation of the SPA}

Consider a general oscillatory integral
\begin{equation}
    I = \int_{-\infty}^{\infty} A(z)\,\te^{\ti\Phi(z)}\,\td z.
\end{equation}
Assume there exists a unique stationary point \(z^*\) such that
\begin{equation}
    \Phi'(z^*)=0 \quad \text{and} \quad \Phi''(z^*)\neq0.
\end{equation}
The main idea of the SPA is that if the phase oscillates rapidly away from \(z^*\) while the amplitude \(A(z)\) varies slowly, the integral is dominated by the region near \(z^*\).

\paragraph{Change of Variables.}  
Introduce a new variable \(u\) defined by
\begin{equation}
    u^2 = -\ti\Bigl(\Phi(z)-\Phi(z^*)\Bigr),
    \label{eq:u_def}
\end{equation}
so that \(u=0\) corresponds to \(z=z^*\). Then, without approximation, the integral can be written as
\begin{equation}
    I = \te^{\ti\Phi(z^*)}\int_{0}^{\infty} A(z(u))\frac{\td z}{\td u}\,\te^{-u^2}\,\td u.
\end{equation}

\paragraph{Expansion Around \(z^*\).}  
Expanding the phase in a Taylor series around \(z=z^*\) gives
\begin{equation}
    \Phi(z) \approx \Phi(z^*) + \frac{1}{2}\Phi''(z^*)(z-z^*)^2 + O\big((z-z^*)^4\big).
\end{equation}
Then, from \eqref{eq:u_def},
\begin{equation}
    u^2 \approx -\ti\,\frac{\Phi''(z^*)}{2}(z-z^*)^2.
\end{equation}
Inverting this relation yields
\begin{equation}
    z \approx z^* + \sqrt{\frac{2}{\Phi''(z^*)}}\,
    \te^{\ti\frac{\pi}{4}\,\sgn\big(\Phi''(z^*)\big)}\,u + O(u^2).
\end{equation}

Similarly, the amplitude is expanded in a Maclaurin series:
\begin{equation}
    A_u(u) \coloneqq A(z(u))\frac{\td z}{\td u} \approx \sum_{k=0}^{\infty}\frac{A_u^{(k)}(0)}{k!}\,u^k.
\end{equation}
Thus, the integral becomes
\begin{equation}
    I = \te^{\ti\Phi(z^*)}\sum_{k=0}^{\infty}\frac{A_u^{(k)}(0)}{k!}\int_{0}^{\infty} u^k\,\te^{-u^2}\,\td u.
\end{equation}
Using the Gamma–function identity
\begin{equation}
    \int_{0}^{\infty}u^k\,\te^{-u^2}\,\td u=\frac{1}{2}\,\Gamma\Bigl(\frac{k+1}{2}\Bigr),
\end{equation}
we obtain
\begin{equation}
    I = \te^{\ti\Phi(z^*)}\sum_{k=0}^{\infty}\frac{A_u^{(k)}(0)}{k!}\,\frac{1}{2}\,\Gamma\Bigl(\frac{k+1}{2}\Bigr).
\end{equation}
It can be shown that, under suitable symmetry conditions, the odd–order terms vanish so that the leading contribution comes from the zeroth–order term.

\paragraph{Leading–Order (Zeroth–Order) Term.}  
At \(u=0\), one finds
\begin{equation}
    A_u(0)=\sqrt{2}\,\frac{A(z^*)}{|\Phi''(z^*)|^{1/2}}\,
    \te^{\ti\frac{\pi}{4}\,\sgn\big(\Phi''(z^*)\big)}.
\end{equation}
Hence, the leading–order approximation is
\begin{equation}
    I_0 = \sqrt{\frac{2\pi}{|\Phi''(z^*)|}}\,
    \te^{\ti\Bigl[\Phi(z^*)+\frac{\pi}{4}\,\sgn\big(\Phi''(z^*)\big)\Bigr]}A(z^*).
\end{equation}

\paragraph{Second–Order Correction.}  
The second–order term involves the second derivative of \(A_u(u)\):
\begin{equation}
    A_u^{(2)}(u)=\frac{\td^2}{\td u^2}\Bigl[A(z(u))\frac{\td z}{\td u}\Bigr]
    =A''(z(u))(z'(u))^3+3A'(z(u))z'(u)z''(u)+A(z(u))z'''(u).
\end{equation}
After expressing the derivatives in terms of \(z\) (and evaluating at \(u=0\)) one obtains (after some algebra):
\begin{multline}
    A_u^{(2)}(0)=2\sqrt{2}\,\frac{\te^{3\ti\frac{\pi}{4}\,\sgn\big(\Phi''(z^*)\big)}}{\Phi''(z^*)\,|\Phi''(z^*)|^{3/2}} \times \\
    \Biggl[A''(z^*)\,\Phi''(z^*)^2-A'(z^*)\,\Phi'''(z^*)\,\Phi''(z^*)
    +\frac{A(z^*)}{12}\Bigl(5\,\Phi'''(z^*)^2-3\,\Phi''''(z^*)\,\Phi''(z^*)\Bigr)\Biggr].
\end{multline}
Thus, the second–order contribution is given by
\begin{multline}
    I_2=\frac{\sqrt{2\pi}}{24}\,\te^{\ti\Phi(z^*)}\,
    \frac{\te^{3\ti\frac{\pi}{4}\,\sgn\big(\Phi''(z^*)\big)}}{\Phi''(z^*)^2\,|\Phi''(z^*)|^{3/2}} \times \\
    \Biggl[A''(z^*)\,\Phi''(z^*)^2-A'(z^*)\,\Phi'''(z^*)\,\Phi''(z^*)
    +\frac{A(z^*)}{12}\Bigl(5\,\Phi'''(z^*)^2-3\,\Phi''''(z^*)\,\Phi''(z^*)\Bigr)\Biggr].
\end{multline}

\paragraph{Error Estimation.}  
A practical measure for the validity of the SPA is to require that
\begin{equation}
    \frac{|I_2|}{|I_0|}\ll1.
    \label{eq:SPAvalidity}
\end{equation}
After some algebra, this condition can be expressed (in a simplified form) as
\begin{equation}
    \frac{A''(z^*)}{A(z^*)\,|\Phi''(z^*)|}
    -\frac{1}{4}\,\frac{|\Phi''''(z^*)|}{|\Phi''(z^*)|^2}
    \ll1.
\end{equation}
In many practical situations (especially when the amplitude is slowly varying and the phase is even), further simplifications are possible.

\subsection{Application to the Rayleigh Integral}

Let us now apply the SPA to the Rayleigh integral in \eqref{eq:RayleighPolar}. Assume that the evaluation is carried out in the horizontal plane (\(z=0\)) and that the target field propagates horizontally (i.e. \(k^P_z(x,y,0)=0\)). Consequently, the stationary point in the vertical direction is \(z^*=0\), independent of \(x\). Moreover, if the source distribution is located in the horizontal plane, the amplitude and phase functions are even in \(z\) so that
\begin{equation}
    A'(x,y,0)=0 \quad \text{and} \quad \Phi'''(x,y,0)=0.
\end{equation}
Similar properties hold for the Green’s function. Then, for the combined amplitude and phase in \eqref{eq:RayleighPolar} (omitting constant multipliers) we have
\begin{equation}
    A'(x,y,0)=0,\quad \Phi'''(x,y,0)=0.
\end{equation}
Under these conditions, the validity condition \eqref{eq:SPAvalidity} simplifies to
\begin{equation}
    \frac{|I_2|}{|I_0|}\approx
    \frac{A''(0)}{A(0)|\Phi''(0)|}
    -\frac{1}{4}\,\frac{|\Phi''''(0)|}{|\Phi''(0)|^2}
    \ll1.
    \label{eq:RayleighSPAvalidity}
\end{equation}
Expressing the phase in the form
\[
\Phi(\vx)=\frac{\omega}{c}\,\hat{\Phi}(\vx),
\]
the condition can be rewritten as
\begin{equation}
    \frac{A''(0)}{A(0)|\hat{\Phi}''(0)|}
    -\frac{1}{4}\,\frac{|\hat{\Phi}''''(0)|}{|\hat{\Phi}''(0)|^2}
    \ll\frac{\omega}{c}.
\end{equation}
When the SPA is used to reduce the \(xz\)–plane integral to a line integral along \(x\), this criterion gives an estimate of the lower–frequency limit (per SSD element). The resulting line integral reads
\begin{multline}
    P(x,y,\omega)\approx2\int_{-\infty}^{\infty}\sqrt{\frac{2\pi}{\ti\,|\Phi^P(\vxo)+\Phi^G(\vx-\vxo)|}} \\
    \times\,\ti\,k_y(x_0,y_0,0)\,P(x_0,y_0,0)\,G(x-x_0,y-y_0,0)\,\td x_0.
\end{multline}
One may then define the overall lower–frequency limit of the approximation as the cutoff frequency for the horizontal (stationary) SSD element \(x_0^*\):
\begin{equation}
    \omega^P(\vx)=c\left[\frac{A''(x_0^*,y_0,0)}{A(x_0^*,y_0,0)|\hat{\Phi}''(x_0^*,y_0,0)|}
    -\frac{1}{4}\,\frac{|\hat{\Phi}''''(x_0^*,y_0,0)|}{|\hat{\Phi}''(x_0^*,y_0,0)|^2}\right],
\end{equation}
where the dependence of \(x_0^*\) on \(\vx\) and the target field is implicit.

\section{Application to a Virtual Point Source}

% (This section will detail the application of the above theory to the case of a virtual point source. Additional derivations and examples can be added here.)

\end{document}
